% ------------------------------------------------------------%
% 2025-11-12
% 2025 - Emerson Ribeiro de Mello - mello@ifsc.edu.br
% ------------------------------------------------------------%

%---------------------------------------------------------%
% ATENÇÃO: O pacote ifscutils é necessário para essa classe
% Para evitar dependências de soluções externas para localizar
% tal arquivo, deve-se incluir a linha abaixo dentro do arquivo
% .tex, já apontando para o caminho relativo do arquivo ifscutils.sty
% \RequirePackage{ifscutils}
%---------------------------------------------------------%

\documentclass[11pt]{ifscarticle}
\RequirePackage{./ifscutils}

\usepackage{xargs}
\usepackage{listings}

\begin{document}

\cabecalhoCurso{Análise e Desenvolvimento de Sistemas}{Segurança da Informação}{Prof. Emerson Ribeiro de Mello}{\url{mello@ifsc.edu.br}}{\textbf{Prática:} OWASP WebGoat e WebWolf}{18/11/2025}

\noindent\textbf{Integrantes:}\newline
Giuliano Marcus Bianco\newline
Luís Renato Freitas de Almeida\newline
Nicolas Arthur Raulino Oliveira

\vspace{0.5cm}

\section{Objetivos}

\begin{itemize}
    \item Configurar e iniciar o ambiente da prática usando Docker Compose
    \item Acessar as aplicações WebGoat e WebWolf
    \item Usar o OWASP ZAP para interceptar e analisar requisições HTTP
    \item Atacar algumas vulnerabilidades do WebGoat
\end{itemize}


\section{Configuração do Ambiente}

O ambiente de laboratório é composto por três aplicações principais:

\begin{itemize}
    \item \textbf{WebGoat}: Aplicação web deliberadamente insegura que contém vulnerabilidades comumente encontradas em aplicações Java. É utilizada para aprender e praticar técnicas de segurança web através de exercícios práticos.
    \item \textbf{WebWolf}: Aplicação complementar que simula o ambiente do atacante. Permite receber e-mails, hospedar arquivos e capturar requisições enviadas pelo WebGoat durante os exercícios.
    \item \textbf{OWASP ZAP}: Proxy de segurança para interceptar e analisar requisições HTTP/HTTPS, permitindo identificar vulnerabilidades e testar ataques.
\end{itemize}

\subsection{Arquivo docker-compose.yml}

O arquivo de configuração do Docker Compose é apresentado no \autoref{docker-compose}.

\lstinputlisting[style=yaml,caption={Configuração do Docker Compose},label={docker-compose}]{../docker-compose.yml}

\section{Iniciando o Ambiente}

\subsection{Iniciar as aplicações}

Para iniciar o WebGoat, WebWolf e OWASP ZAP, execute o comando abaixo:

\begin{lstlisting}[style=shell]
docker-compose up --build -d
\end{lstlisting}

\section{Acessando as Aplicações}

\subsection{Acessando o OWASP ZAP}

O OWASP ZAP possui interface gráfica que será exibida automaticamente após a inicialização. O ZAP atua como um \textit{Man-in-the-Middle} (MITM), posicionando-se entre o navegador e o servidor web para interceptar, analisar e modificar o tráfego HTTP/HTTPS. A interface permite visualizar todas as requisições e respostas capturadas.

\subsection{Acessando o WebGoat}

O WebGoat é a aplicação principal desta prática. Ela contém diversas vulnerabilidades comumente encontradas em aplicações Java. Para acessar o WebGoat:

\begin{enumerate}
    \item Abra a interface gráfica do OWASP ZAP
    \item Use o navegador integrado do ZAP (disponível na interface)
    \item No navegador integrado, acesse:
    \begin{lstlisting}[style=shell]
http://webgoat:8080/WebGoat\end{lstlisting}
\end{enumerate}

O navegador integrado do ZAP já está configurado para interceptar automaticamente todo o tráfego HTTP/HTTPS, permitindo analisar as requisições e respostas enviadas pelo WebGoat.

\begin{important}
    Use \texttt{webgoat:8080} ao invés de \texttt{localhost:8080} no navegador integrado do ZAP, pois as aplicações estão na mesma rede interna e se comunicam pelo nome do serviço.
\end{important}

\subsection{Acessando o WebWolf}

O WebWolf simula o ambiente do atacante e é utilizado durante os exercícios para receber dados enviados pelo WebGoat. Diferente do WebGoat, o WebWolf deve ser acessado através do navegador na sua máquina local:

\begin{enumerate}
    \item Abra seu navegador web
    \item Acesse:
    \begin{lstlisting}[style=shell]
http://localhost:9090/WebWolf
\end{lstlisting}
\end{enumerate}

O WebWolf permite visualizar e-mails recebidos, arquivos hospedados e requisições capturadas durante os exercícios práticos do WebGoat.

\vspace*{\fill}
\licenciamentoLivre

\end{document}
